\documentclass[11pt]{article}
\usepackage[margin=0.25in]{geometry}
\usepackage{amsmath}
\usepackage{amssymb}
\usepackage{listings}
\usepackage{xcolor}
\usepackage{hyperref}
\usepackage{graphicx}
\usepackage{enumitem}

% Code listing style
\lstdefinestyle{pythonstyle}{
    language=Python,
    basicstyle=\ttfamily\small,
    keywordstyle=\color{blue}\bfseries,
    commentstyle=\color{green!50!black},
    stringstyle=\color{red},
    numbers=left,
    numberstyle=\tiny\color{gray},
    stepnumber=1,
    numbersep=5pt,
    backgroundcolor=\color{gray!10},
    showspaces=false,
    showstringspaces=false,
    showtabs=false,
    frame=single,
    tabsize=2,
    captionpos=b,
    breaklines=true,
    breakatwhitespace=false,
    escapeinside={(*@}{@*)},
    morekeywords={True,False,None,and,or,not,in,is}
}

\lstset{style=pythonstyle}

\title{\textbf{Step 6: Forming Polygons from Edges} \\
\large Algorithm Explanation and Implementation}
\author{3D Solid Reconstruction Pipeline}
\date{October 2025}

\begin{document}

\maketitle

\begin{abstract}
This document provides a comprehensive explanation of the polygon formation algorithm (Step 6) in the 3D solid reconstruction pipeline. The algorithm takes a set of edges lying on a planar face and constructs the face's boundary polygon plus any interior holes, handling complex cases of overlapping and touching polygons through geometric decomposition and iterative merging.
\end{abstract}

\tableofcontents
\newpage

\section{High-Level Overview}

Step 6 transforms a set of edges lying on a planar face into a structured face representation consisting of:
\begin{itemize}
    \item An \textbf{outer boundary polygon} (largest area)
    \item Zero or more \textbf{interior holes} (polygons fully contained inside boundary)
    \item \textbf{Alternate boundary candidates} (for topological ambiguity resolution)
\end{itemize}

The algorithm handles complex scenarios including:
\begin{itemize}
    \item Multiple overlapping polygons requiring decomposition
    \item Touching polygons that share edges but represent separate faces
    \item Colinear vertex expansion (intermediate vertices on edges)
    \item Disconnected components requiring multiple iterations
\end{itemize}

\section{Core Algorithm Flow}

\begin{lstlisting}[caption={High-level algorithm structure},label={lst:highlevel}]
For each face:
  WHILE (unused vertices remain):
    ITERATION:
      1. Find ALL possible cycles from remaining edges
      2. Choose largest cycle as boundary candidate
      3. Find other polygons that share edges with boundary
      4. MERGE LOOP: For each polygon sharing edges:
         a. Classify relationship (touching, overlapping, subset)
         b. Apply appropriate merge strategy
         c. Update boundary or add to alternate list
      5. Classify remaining polygons (inside=holes, outside=next iter)
      6. Store boundary + holes as one face
\end{lstlisting}

\section{Detailed Algorithm Components}

\subsection{Initial Setup and 2D Projection}

\begin{lstlisting}[caption={Initialization and projection},label={lst:init}]
# Project face vertices to 2D coordinate system
if abs(normal[2]) < 0.9:
    u = cross(normal, [0, 0, 1])
else:
    u = cross(normal, [1, 0, 0])
u = u / norm(u)
v = cross(normal, u)

verts_2d = {}
for v_idx in vertices_on_face:
    vert_3d = selected_verts[v_idx]
    verts_2d[v_idx] = (dot(vert_3d, u), dot(vert_3d, v))

# Build adjacency list from edges
adjacency = {}
for edge in edge_set:
    v1, v2 = edge
    adjacency[v1].append(v2)
    adjacency[v2].append(v1)
\end{lstlisting}

\textbf{Purpose:} Transform 3D problem into 2D polygon problem using orthonormal basis $(u, v)$ in the face plane.

\subsection{Outer Iteration Loop}

\begin{lstlisting}[caption={Multi-iteration structure for disconnected components},label={lst:outerloop}]
iteration = 1
while len(used_vertices) < len(vertices_on_face):
    remaining_edges = edge_set - used_edges
    remaining_verts = vertices with edges in remaining_edges
    
    # Process one connected component per iteration
    ...
\end{lstlisting}

\textbf{Purpose:} Handle disconnected components. Each iteration processes one connected group of edges.

\textbf{Example:}
\begin{itemize}
    \item Iteration 1: Main boundary (8 vertices)
    \item Iteration 2: Interior hole (4 vertices)
\end{itemize}

\subsection{Comprehensive Cycle Discovery}

\subsubsection{DFS-Based Algorithm}

\begin{lstlisting}[caption={Find all cycles using depth-first search with backtracking},label={lst:cyclediscovery}]
def find_all_cycles_from_edges(edges, verts_2d):
    # Build adjacency list
    adjacency = build_adjacency(edges)
    all_polygons = []
    found_edge_sets = set()  # Frozenset for duplicate detection
    
    def dfs_find_cycles(start_vertex, current, path, path_edges):
        """DFS with backtracking to find all cycles."""
        for neighbor in adjacency[current]:
            edge = (min(current, neighbor), max(current, neighbor))
            
            if edge not in edge_set or edge in path_edges:
                continue
            
            # Check if we can close the cycle
            if neighbor == start_vertex and len(path) >= 3:
                cycle_edges = path_edges | {edge}
                if frozenset(cycle_edges) not in found_edge_sets:
                    found_edge_sets.add(frozenset(cycle_edges))
                    normalized = normalize_polygon(path)
                    all_polygons.append(normalized)
                continue
            
            # Continue DFS if neighbor not in path
            if neighbor not in path:
                dfs_find_cycles(start_vertex, neighbor,
                               path + [neighbor], path_edges | {edge})
    
    # Start DFS from each vertex-neighbor pair
    for start_vertex in vertices:
        for first_neighbor in adjacency[start_vertex]:
            dfs_find_cycles(start_vertex, first_neighbor,
                           [start_vertex, first_neighbor], {edge})
    
    return all_polygons
\end{lstlisting}

\textbf{Key Innovation:} Uses DFS with backtracking to explore \textit{all} possible paths, ensuring no valid cycles are missed.

\textbf{Example (Face 16, seed 255):}
\begin{itemize}
    \item \textbf{Before:} Greedy approach found 2 polygons, chose 4-vertex (area=326.03)
    \item \textbf{After:} DFS found 3 polygons, chose 8-vertex (area=617.86) — 89\% larger!
\end{itemize}

\subsubsection{Boundary Selection}

\begin{lstlisting}[caption={Selecting initial boundary candidate},label={lst:boundaryselect}]
# Calculate areas and sort
polys_with_area = []
for poly in all_possible_polygons:
    poly_shapely = Polygon([verts_2d[v] for v in poly])
    if poly_shapely.is_valid and poly_shapely.area > 1e-10:
        polys_with_area.append({
            'vertices': poly,
            'shapely': poly_shapely,
            'area': poly_shapely.area
        })

# Sort by area (largest first)
polys_with_area.sort(
    key=lambda p: (p['area'], len(p['vertices'])),
    reverse=True
)

# Choose largest as initial boundary
boundary_poly = polys_with_area[0]['vertices']
all_other_polygons = [p['vertices'] for p in polys_with_area[1:]]
\end{lstlisting}

\subsection{The Merge Loop: Polygon Relationship Classification}

This is the heart of the algorithm, where polygons sharing edges with the boundary are processed through iterative merging or decomposition.

\begin{lstlisting}[caption={Main merge loop structure},label={lst:mergeloop}]
merged_any = True
merge_count = 0

while merged_any and merge_count < 20:
    merged_any = False
    merge_count += 1
    
    boundary_edges = get_polygon_edges(boundary_poly)
    
    for poly in all_other_polygons[:]:
        poly_edges = get_polygon_edges(poly)
        shared_edges = boundary_edges (*@$\cap$@*) poly_edges
        
        if not shared_edges:
            continue
        
        # Classify and process relationship
        ...
\end{lstlisting}

\subsubsection{Type 1: Full Overlap (Subset)}

\begin{lstlisting}[caption={Handling complete overlap},label={lst:fullover}]
poly_unique_edges = poly_edges - shared_edges

if not poly_unique_edges:
    # Poly has no unique edges, fully overlaps
    used_edges.update(shared_edges)
    all_other_polygons.remove(poly)
    merged_any = True
    break
\end{lstlisting}

\textbf{Condition:} Polygon has no edges outside the boundary.

\textbf{Action:} Mark shared edges as used, remove polygon.

\subsubsection{Type 2: Colinear Expansion}

\begin{lstlisting}[caption={Colinear vertex insertion},label={lst:colinear}]
# All unique vertices lie on boundary AND polygon area ~0
if boundary_count == len(unique_verts) and poly_area < 1e-6:
    # Poly inserts intermediate vertices along boundary edge
    boundary_poly = poly  # Replace with expanded version
    boundary_2d = [verts_2d[v] for v in boundary_poly]
    boundary_shapely = Polygon(boundary_2d)
    
    used_edges.update(shared_edges)
    all_other_polygons.remove(poly)
    merged_any = True
\end{lstlisting}

\textbf{Condition:} Polygon has zero area and only adds vertices along existing boundary edges.

\textbf{Action:} Replace boundary with expanded polygon (more vertices, same shape).

\subsubsection{Type 3: Separate Touching Faces}

\begin{lstlisting}[caption={Two faces sharing an edge},label={lst:touching}]
# Remove shared edges from BOTH polygons
boundary_unique_edges = boundary_edges - shared_edges
poly_unique_edges = poly_edges - shared_edges

# Try to form polygons from remaining edges
boundary_remaining = find_all_cycles(boundary_unique_edges)
poly_remaining = find_all_cycles(poly_unique_edges)

if boundary_remaining and poly_remaining:
    # Both form valid polygons -> SEPARATE FACES
    
    # Choose larger as new boundary
    if area(poly_remaining) > area(boundary_remaining):
        swap(boundary_poly, poly)
        all_other_polygons.append(old_boundary)
    
    # Mark shared edge as boundary between faces
    used_edges.update(shared_edges)
    all_other_polygons.remove(poly)
    merged_any = True
\end{lstlisting}

\textbf{Condition:} After removing shared edges, both polygons still form valid cycles.

\textbf{Example:}
\begin{align*}
\text{Face A:} &\quad [1,2,3,4] \\
\text{Face B:} &\quad [3,4,5,6] \\
\text{Shared:} &\quad (3,4) \\
\text{Result:} &\quad \text{Keep both as separate faces}
\end{align*}

\subsubsection{Type 4: Overlapping Polygons (Geometric Decomposition)}

\textbf{Condition:} One or both polygons don't form valid cycles after removing shared edges.

\textbf{Key Difference from Type 3:} Unlike Type 3 where both polygons remain valid after edge removal (separate touching faces), here we have true overlap that requires decomposition.

When polygons genuinely overlap, use Shapely geometric operations to decompose them:

\begin{lstlisting}[caption={Geometric decomposition of overlapping polygons},label={lst:geometric}]
# Convert to Shapely polygons
boundary_shapely = Polygon([verts_2d[v] for v in boundary_poly])
poly_shapely = Polygon([verts_2d[v] for v in poly])

# Geometric operations decompose into 3 possible regions:
intersection = boundary_shapely.intersection(poly_shapely)     # Shared area
diff_boundary = boundary_shapely.difference(poly_shapely)      # Boundary-only
diff_poly = poly_shapely.difference(boundary_shapely)          # Poly-only

# Collect non-empty regions (area > threshold)
result_regions = []
if intersection.area > 1e-6:
    result_regions.append(('intersection', intersection))
if diff_boundary.area > 1e-6:
    result_regions.append(('boundary_diff', diff_boundary))
if diff_poly.area > 1e-6:
    result_regions.append(('poly_diff', diff_poly))

# Convert Shapely polygons back to vertex indices
converted_polys = []
for region_name, shapely_geom in result_regions:
    coords = list(shapely_geom.exterior.coords[:-1])
    matched_verts = match_coords_to_vertices(coords, verts_2d)
    converted_polys.append((region_name, matched_verts, shapely_geom.area))

# Sort by area (largest first)
converted_polys.sort(key=lambda x: x[2], reverse=True)

# CRITICAL: Original boundary is DISCARDED and REPLACED
boundary_poly = converted_polys[0][1]  # Largest decomposed region

# Add remaining decomposed regions back to list (with duplicate check)
for region_name, region_verts, area in converted_polys[1:]:
    if not is_duplicate(region_verts, all_other_polygons):
        all_other_polygons.append(region_verts)
        
# Original polygons (boundary and poly) are now gone,
# replaced by their geometric decomposition
\end{lstlisting}

\textbf{Important Notes:}
\begin{enumerate}
    \item The \textbf{original boundary polygon is NOT retained} — it's completely replaced by the decomposed regions
    \item We may get 1, 2, or 3 regions depending on overlap configuration:
    \begin{itemize}
        \item Complete overlap: Only intersection
        \item Partial overlap: Intersection + 1 or 2 differences
        \item Edge-touching: Only differences (intersection has zero area)
    \end{itemize}
    \item The largest decomposed region becomes the new boundary
    \item Smaller regions are added back for further processing
\end{enumerate}

\textbf{Example (Face 10, seed 55):}

\begin{table}[h]
\centering
\begin{tabular}{|l|l|l|}
\hline
\textbf{Polygon} & \textbf{Vertices} & \textbf{Area} \\ \hline
Boundary (initial) & 10 vertices & 996 \\ \hline
Overlapping poly & 12 vertices & 964 \\ \hline
Shared edges & \multicolumn{2}{c|}{9 edges} \\ \hline
\hline
\multicolumn{3}{|c|}{\textit{After removing shared edges:}} \\ \hline
Boundary remaining & 1 edge & No valid cycle \\ \hline
Poly remaining & 3 edges & No valid cycle \\ \hline
\hline
\multicolumn{3}{|c|}{\textit{Geometric decomposition:}} \\ \hline
Intersection & 10 vertices & 996 \\ \hline
Diff\_poly & 4 vertices & 4.6 \\ \hline
\hline
\textbf{Final boundary} & \textbf{12 vertices} & \textbf{964} \\ \hline
\end{tabular}
\caption{Geometric decomposition example - original 10-vert and 12-vert polygons are discarded and replaced}
\end{table}

\textbf{Result:} The original 10-vertex boundary is discarded. The intersection (10 verts, area=996) becomes the new boundary, and diff\_poly (4 verts) is added back to the list.

\subsubsection{Type 3 vs Type 4 Comparison}

To clarify the critical difference between separate touching faces (Type 3) and overlapping polygons (Type 4):

\begin{table}[h]
\centering
\begin{tabular}{|p{3cm}|p{5cm}|p{5cm}|}
\hline
\textbf{Aspect} & \textbf{Type 3: Separate Touching} & \textbf{Type 4: Overlapping} \\ \hline
\hline
\textbf{Condition} & Both form valid cycles after removing shared edges & One or both DON'T form valid cycles \\ \hline
\hline
\textbf{Topology} & Two distinct faces sharing a boundary edge & Polygons genuinely overlap (shared interior area) \\ \hline
\hline
\textbf{Approach} & Keep both polygons separately & Geometric decomposition (Shapely) \\ \hline
\hline
\textbf{Original polygons} & \textbf{Both retained} (larger becomes boundary, smaller processed next iteration) & \textbf{Both discarded}, replaced by decomposed regions \\ \hline
\hline
\textbf{Shared edges} & Marked as used (boundary between faces) & Edges absorbed into decomposed regions \\ \hline
\hline
\textbf{Example} & Two room faces sharing a wall & Overlapping circle and square \\ \hline
\end{tabular}
\caption{Type 3 vs Type 4 classification}
\end{table}

\textbf{Visual Example:}

\begin{verbatim}
Type 3 (Touching):
  +-------+-------+
  | Face  | Face  |    After removing shared edge (3,4):
  |   A   |   B   |    Face A: [1,2,3,4] -> still valid
  +-------+-------+    Face B: [3,4,5,6] -> still valid
  Result: Keep both faces

Type 4 (Overlapping):
  +-------+           After removing shared edges:
  | Bndry |           Boundary: 1 edge -> NO cycle
  |   +---+---+       Poly: 3 edges -> NO cycle
  |   | Overlap |     
  +---+---+     |     Must use geometric decomposition:
      | Poly  |       intersection + differences
      +-------+       
  Result: Decompose into regions, discard originals
\end{verbatim}

\subsection{Duplicate Detection (Critical Fix)}

\begin{lstlisting}[caption={Preventing infinite loops through duplicate detection},label={lst:duplicate}]
def polygons_are_same(poly1, poly2):
    """Check if two polygons are the same (allowing rotation/reversal)."""
    if len(poly1) != len(poly2):
        return False
    
    # Normalize (start from minimum vertex)
    def normalize(poly):
        min_idx = poly.index(min(poly))
        return poly[min_idx:] + poly[:min_idx]
    
    norm1 = normalize(poly1)
    norm2 = normalize(poly2)
    
    # Check forward and reverse directions
    if norm1 == norm2:
        return True
    norm2_rev = [norm2[0]] + norm2[1:][::-1]
    return norm1 == norm2_rev

# Use in merge loop
for region_name, region_verts, area in converted_polys[1:]:
    is_duplicate = False
    for existing_poly in all_other_polygons:
        if polygons_are_same(region_verts, existing_poly):
            is_duplicate = True
            break
    
    if not is_duplicate:
        all_other_polygons.append(region_verts)
\end{lstlisting}

\textbf{Problem Solved:} Before this fix, rotated versions of the same polygon were repeatedly added, causing infinite loops.

\textbf{Example (Face 16 before fix):}
\begin{itemize}
    \item Iteration 1: Add [105, 30, 8, 49, 15, 59, 64]
    \item Iteration 2: Add [30, 8, 49, 15, 59, 64, 105] $\leftarrow$ Duplicate (rotated)
    \item Iteration 3: Add [8, 49, 15, 59, 64, 105, 30] $\leftarrow$ Duplicate (rotated)
    \item ... (continues forever)
\end{itemize}

\textbf{After fix:} Detects rotation, skips duplicate $\rightarrow$ loop terminates in 3 iterations.

\subsection{Final Classification and Storage}

\begin{lstlisting}[caption={Classifying remaining polygons as holes or external},label={lst:classify}]
# Mark boundary vertices as used
used_vertices.update(boundary_poly)

# Classify remaining polygons
holes = []
outside_polys = []

for poly in all_other_polygons:
    # Check if all vertices are inside boundary
    all_inside = True
    for v in poly:
        if v not in boundary_poly:
            point = Point(verts_2d[v])
            if not boundary_shapely.contains(point):
                all_inside = False
                break
    
    if all_inside:
        holes.append(poly)
    else:
        outside_polys.append(poly)

# Store face (boundary + holes)
faces.append({
    'boundary': boundary_poly,
    'holes': holes,
    'alternates': alternates_candidates
})

# Outside polygons will be processed in next iteration
\end{lstlisting}

\textbf{Classification rules:}
\begin{itemize}
    \item \textbf{Hole:} All vertices lie inside boundary
    \item \textbf{External:} At least one vertex outside boundary
\end{itemize}

\section{Algorithm Complexity Analysis}

\subsection{Time Complexity}

\begin{align*}
T(V, E, C, M) &= O(V \times E \times C \times M) \\
\text{where:} \quad
V &= \text{number of vertices} \\
E &= \text{number of edges} \\
C &= \text{number of cycles per face (typically small)} \\
M &= \text{merge iterations (bounded to 20)}
\end{align*}

\textbf{Breakdown:}
\begin{itemize}
    \item Cycle discovery: $O(V \times E \times C)$ — DFS explores all paths
    \item Merge loop: $O(M \times P \times E)$ where $P$ = polygons per iteration
    \item Geometric operations: $O(V)$ per Shapely operation
\end{itemize}

\subsection{Space Complexity}

\begin{equation}
S(E, P) = O(E + P \times V_{avg})
\end{equation}

\begin{itemize}
    \item Edge storage: $O(E)$
    \item Polygon storage: $O(P \times V_{avg})$ where $V_{avg}$ = average vertices per polygon
    \item Adjacency list: $O(V + E)$
\end{itemize}

\section{Real-World Test Cases}

\subsection{Simple Face (4 edges, 4 vertices)}

\begin{lstlisting}[caption={Simple rectangular face},label={lst:simple}]
Edges: [(1,2), (2,3), (3,4), (4,1)]
Vertices: [1, 2, 3, 4]

Iteration 1:
  - Find 1 cycle: [1,2,3,4]
  - No other polygons
  Result: Boundary=[1,2,3,4], no holes
\end{lstlisting}

\subsection{Face with Interior Hole}

\begin{lstlisting}[caption={Face with rectangular hole},label={lst:hole}]
Edges: 12 edges forming outer and inner rectangles
Vertices: 8 vertices (4 outer, 4 inner)

Iteration 1:
  - Find 2 cycles: 
      outer [1,2,3,4,5,6,7,8]
      inner [3,4,9,10]
  - Choose outer as boundary (larger area)
  - Classify inner: all vertices INSIDE -> HOLE
  Result: Boundary=[1,2,3,4,5,6,7,8], Hole=[3,4,9,10]
\end{lstlisting}

\subsection{Complex Overlapping (Face 10, seed 55)}

\begin{lstlisting}[caption={Complex face with multiple overlapping polygons},label={lst:complex}]
Edges: 40 edges
Vertices: 15 vertices

Iteration 1:
  - Find 35 cycles from 40 edges
  - Choose 10-vert as boundary (largest area)
  - Process 34 other polygons through merge loop
  - Merge iterations: 20 (decompose overlaps)
  - Geometric operations: 15+ decompositions
  Final: Boundary=12-vert (merged from multiple regions)
         Area=964.73, includes all edge groups
\end{lstlisting}

\section{Performance Metrics}

\subsection{Seed 55 Results}

\begin{table}[h]
\centering
\begin{tabular}{|l|c|}
\hline
\textbf{Metric} & \textbf{Value} \\ \hline
Total faces & 31 \\ \hline
Faces with holes & 3 \\ \hline
Face 7 vertices & 16 (with 1 hole) \\ \hline
Face 10 vertices & 12 (corrected from 10) \\ \hline
Merge iterations (Face 10) & 20 \\ \hline
Solid validity & Valid \\ \hline
Reconstruction time & $<$ 5 seconds \\ \hline
\end{tabular}
\caption{Seed 55 reconstruction metrics}
\end{table}

\subsection{Seed 255 Results}

\begin{table}[h]
\centering
\begin{tabular}{|l|c|c|}
\hline
\textbf{Metric} & \textbf{Before Fix} & \textbf{After Fix} \\ \hline
Face 16 cycles found & 2 & 3 \\ \hline
Face 16 boundary vertices & 4 & 8 (initially) \\ \hline
Face 16 boundary area & 326.03 & 617.86 \\ \hline
Face 16 merge iterations & 20+ (infinite) & 3 \\ \hline
Unused edges & 6 & 1 \\ \hline
\end{tabular}
\caption{Seed 255 Face 16 improvements}
\end{table}

\section{Key Algorithm Properties}

\subsection{Completeness}

\textbf{Theorem:} The DFS-based cycle discovery algorithm finds all simple cycles in the edge graph.

\textbf{Proof sketch:} By exploring all possible paths through backtracking, every valid cycle starting from each vertex is discovered. The frozenset-based duplicate detection ensures each unique cycle is recorded exactly once.

\subsection{Correctness}

\textbf{Theorem:} The geometric decomposition approach correctly handles all polygon overlap configurations.

\textbf{Justification:} Shapely's intersection and difference operations are based on proven computational geometry algorithms (GEOS library). The approach handles:
\begin{itemize}
    \item Partial overlap $\rightarrow$ decompose into intersection + differences
    \item Complete overlap $\rightarrow$ intersection = smaller polygon
    \item Touching $\rightarrow$ intersection has zero area
    \item Disjoint $\rightarrow$ intersection is empty
\end{itemize}

\subsection{Robustness}

\begin{enumerate}
    \item \textbf{Floating-point tolerance:} Uses $\epsilon = 10^{-6}$ for area comparisons
    \item \textbf{Degenerate cases:} Handles zero-area polygons (colinear vertices)
    \item \textbf{Bounded iterations:} Merge loop limited to 20 iterations (prevents runaway)
    \item \textbf{Fallback strategies:} Multiple approaches for different polygon configurations
\end{enumerate}

\section{Comparison with Previous Approach}

\subsection{Old Algorithm (Greedy Path Exploration)}

\begin{lstlisting}[caption={Old greedy approach},label={lst:oldapproach}]
# Find cycles using greedy single-path building
for start_vertex in vertices:
    for first_neighbor in adjacency[start_vertex]:
        path = [start_vertex, first_neighbor]
        current = first_neighbor
        
        while True:
            neighbors = [n for n in adjacency[current] if n != prev]
            for neighbor in neighbors:
                if valid_edge(current, neighbor):
                    path.append(neighbor)
                    current = neighbor
                    break  # (*@\textbf{COMMITS to first valid neighbor!}@*)
            
            if can_close_cycle():
                polygons.append(path)
                break
\end{lstlisting}

\textbf{Problems:}
\begin{enumerate}
    \item \textbf{Incomplete:} Misses cycles requiring different neighbor choices
    \item \textbf{No backtracking:} Once committed to a path, never explores alternatives
    \item \textbf{Suboptimal:} Often selects smaller polygons over larger ones
\end{enumerate}

\subsection{New Algorithm (DFS with Backtracking)}

\textbf{Improvements:}
\begin{enumerate}
    \item \textbf{Complete exploration:} Tries all valid neighbors at each vertex
    \item \textbf{Backtracking:} Recursively explores alternative paths
    \item \textbf{Optimal selection:} Finds all cycles, then chooses by area
    \item \textbf{Duplicate prevention:} Frozenset edge-set comparison
\end{enumerate}

\subsection{Impact Summary}

\begin{table}[h]
\centering
\begin{tabular}{|l|c|c|c|}
\hline
\textbf{Test Case} & \textbf{Old} & \textbf{New} & \textbf{Improvement} \\ \hline
Face 16 cycles found & 2 & 3 & +50\% \\ \hline
Face 16 boundary area & 326 & 618 & +89\% \\ \hline
Face 16 unused edges & 6 & 1 & -83\% \\ \hline
Merge loop termination & Never & 3 iter & Fixed \\ \hline
Solid validity (seed 55) & Valid & Valid & Maintained \\ \hline
\end{tabular}
\caption{Algorithm improvement metrics}
\end{table}

\section{Implementation Details}

\subsection{Helper Functions}

\subsubsection{Polygon Edge Extraction}

\begin{lstlisting}[caption={Get ordered edges from polygon},label={lst:getedges}]
def get_polygon_edges(poly):
    """Return set of edges in polygon."""
    edges = set()
    for i in range(len(poly)):
        v1 = poly[i]
        v2 = poly[(i + 1) % len(poly)]
        edge = (min(v1, v2), max(v1, v2))
        edges.add(edge)
    return edges
\end{lstlisting}

\subsubsection{Colinear Vertex Expansion}

\begin{lstlisting}[caption={Insert colinear intermediate vertices},label={lst:expand}]
def expand_colinear_edges_in_polygon(poly, edge_set, verts_2d):
    """
    For each edge in polygon, check if there are intermediate
    vertices lying on that edge. Insert them in order.
    """
    expanded = []
    for i in range(len(poly)):
        v1 = poly[i]
        v2 = poly[(i + 1) % len(poly)]
        expanded.append(v1)
        
        # Find all vertices on segment v1->v2
        segment = (verts_2d[v1], verts_2d[v2])
        intermediate = []
        
        for v in edge_set:
            if v != v1 and v != v2:
                if point_on_segment(verts_2d[v], segment):
                    intermediate.append(v)
        
        # Sort by distance from v1
        intermediate.sort(key=lambda v: distance(verts_2d[v1], verts_2d[v]))
        expanded.extend(intermediate)
    
    return expanded
\end{lstlisting}

\subsection{Vertex Matching from Shapely}

\begin{lstlisting}[caption={Convert Shapely polygon back to vertex indices},label={lst:match}]
def match_coords_to_vertices(coords, verts_2d, tolerance=0.01):
    """
    Match Shapely polygon coordinates to original vertex indices.
    """
    matched_verts = []
    
    for coord in coords:
        # Find closest vertex
        min_dist = float('inf')
        best_v = None
        
        for v_idx, v_2d in verts_2d.items():
            dist = sqrt((v_2d[0] - coord[0])**2 + (v_2d[1] - coord[1])**2)
            if dist < min_dist:
                min_dist = dist
                best_v = v_idx
        
        if best_v is not None and min_dist < tolerance:
            matched_verts.append(best_v)
    
    return matched_verts
\end{lstlisting}

\section{Future Enhancements}

\subsection{Potential Optimizations}

\begin{enumerate}
    \item \textbf{Early termination:} Stop cycle discovery when all edges are used
    \item \textbf{Cycle filtering:} Prioritize larger/simpler cycles before exploring complex ones
    \item \textbf{Heuristic pruning:} Skip paths unlikely to yield new cycles
    \item \textbf{Parallel processing:} Independent face processing can be parallelized
\end{enumerate}

\subsection{Alternative Approaches}

\begin{enumerate}
    \item \textbf{Johnson's algorithm:} Standard all-simple-cycles algorithm (more complex but potentially faster for dense graphs)
    \item \textbf{Hierarchical cycle detection:} Find fundamental cycles first, combine later
    \item \textbf{Planarity-aware:} Exploit planar graph properties (faces in planar graphs)
    \item \textbf{Machine learning:} Train model to predict optimal boundary from features
\end{enumerate}

\section{Conclusion}

The polygon formation algorithm (Step 6) successfully addresses the complex problem of reconstructing face boundaries from edge sets through:

\begin{enumerate}
    \item \textbf{Comprehensive cycle discovery} using DFS with backtracking
    \item \textbf{Geometric decomposition} of overlapping polygons via Shapely operations
    \item \textbf{Iterative refinement} through the merge loop
    \item \textbf{Duplicate detection} preventing infinite loops
    \item \textbf{Multi-iteration support} for disconnected components and holes
\end{enumerate}

The algorithm demonstrates strong performance across diverse test cases, correctly handling simple faces, faces with holes, and complex overlapping configurations. The recent improvements (DFS cycle discovery + duplicate detection) have significantly enhanced both correctness and robustness, as evidenced by the Face 16 (seed 255) improvements and maintained validity for seed 55.

This algorithm forms a critical component of the 3D solid reconstruction pipeline, transforming raw edge connectivity into structured face representations suitable for topological solid modeling.

\end{document}
